%%%%%%%%%%  scan idx 15  |  folio 2  %%%%%%%%%%
de
sorte que $T_f$ joue le rôle d'un \underline{fibré tangent relatif virtuel} de $X$ sur $Y$. Dans le cas où le morphisme $f$ est lisse (i.e.\ à application tangente partout surjective), on a simplement
\[
T_f=T_{X/Y}\qquad\text{(fibré tangent le long des fibres)},
\]
tandis que dans le cas où $f:X\to Y$ est une immersion, on trouve
\[
T_f=-\check N_{X/Y},
\]
où $\check N_{X/Y}$ désigne le faisceau normal de $X$ dans $Y$.

Un des buts principaux du Séminaire est de généraliser (1.2) simultanément dans deux directions~:
\begin{enumerate}[label=\alph*)]
\item Se débarasser de l'hypothèse de l'existence d'un corps de base $k$.
\item Remplacer les hypothèses de régularité sur $Y,X$ par une hypothèse de «~régularité locale~» de $f$.
\end{enumerate}
Enfin, chemin faisant nous nous occuperons également du problème~:
\begin{enumerate}[label=\alph*),start=3]
\item Éliminer les hypothèses de quasi-projectivité qui, en l'absence d'un corps de base, s'expriment par l'existence de Modules inversibles amples sur $X$ et sur $Y$.
\end{enumerate}

\subsection*{2.}
Examinons d'abord la généralisation a), en gardant cependant les hypothèses de régularité et d'existence de Modules inversibles amples sur $X$ et sur $Y$.

La définition de $K(X)$, $K(Y)$ et de l'homomorphisme $f_*:K(X)\to K(Y)$ n'offre alors pas de nouveau problème, grâce au fait que $X$ et $Y$ sont réguliers. La voie la plus naturelle pour donner un sens à (1.2) semble donc consister en la définition d'\underline{anneaux de Chow} $A(X),A(Y)$ et d'un homomorphisme de groupes
\[
f_*:A(X)\longrightarrow A(Y),
\]
en l'établissement d'une théorie des classes de Chern, fournissant des applications
\[
c_i:K(X)\longrightarrow A(X),
\]
et de même pour $Y$, et enfin en la description d'un élément fibré tangent relatif virtuel
\[
T_f\in K(X).
\]

